%%%%%%%%%%%%%%%%%%%%%%%%%%%%%%%%%%%%%%%%%%%%%%%%%%%%%%%
% NOTAS TEORIA DE TRANSPORTE OPTIMO
% FCEN-UBA, segundo cuatrimestre de 2026. Fuente unica de las notas
% y de las guias de ejercicios (herramientas/generar_practicas.py).
%%%%%%%%%%%%%%%%%%%%%%%%%%%%%%%%%%%%%%%%%%%%%%%%%%%%%%%

\documentclass[11pt,a4paper]{amsart}
\usepackage{xr-hyper}
\newcounter{chapter}
\usepackage{fancyhdr}
\usepackage{geometry} \geometry{top=2cm,bottom=2.5cm,left=2.5cm,right=2.5cm}

\usepackage{amssymb}
\usepackage{hyperref}
\usepackage{listings}
\usepackage{mathtools}
\usepackage{enumitem}

% Fuzz -------------------------------------------------------------------
% Line spacing -----------------------------------------------------------

%%%%%%%%%fin pegada%%%%%%%%%%%%%%%%%%5

\usepackage[T1]{fontenc}
\usepackage[spanish]{babel}
\usepackage{graphicx}
\usepackage{array}
\usepackage{tikz}

%% Colores del Instituto de Cálculo y del Departamento de Matemática (portada y enlaces)
\definecolor{icnavy}{HTML}{013274}
\definecolor{iccyan}{HTML}{029BDF}
\definecolor{icmag}{HTML}{E2017B}
\definecolor{dmred}{HTML}{A8000C}
\hypersetup{colorlinks=true,linkcolor=icnavy,citecolor=icnavy,urlcolor=icnavy}
\usetikzlibrary{babel}



\parskip 3pt

%%
%% OPERADORES
%%
\renewcommand{\Re}{\operatorname{Re}}
\newcommand{\tr}{\operatorname{tr}}
\newcommand{\id}{\operatorname{id}}
\def\div{\operatorname {\mathrm{div}}}
\def\Dom{\operatorname {\mathrm{Dom}}}
\def\argmin{\operatorname {\mathrm{arg\, min}}}
\def\argmax{\operatorname {\mathrm{arg\, max}}}
\def\gen{\operatorname {\mathrm{gen}}}
\def\Im{\operatorname {\mathrm{Im}}}
\def\Nu{\operatorname {\mathrm{Nu}}}
\def\sop{\operatorname {\mathrm{sop}}}
\def\diam{\operatorname {\mathrm{diam}}}
\def\dist{\operatorname {\mathrm{dist}}}
\def\var{\operatorname {\mathrm{Var}}}
\def\codim{\operatorname {\mathrm{codim}}}
\def\pint{\operatorname {--\!\!\!\!\!\int\!\!\!\!\!--}}

%%
%% EPSILON
%%
\newcommand{\ve}{\varepsilon}

%%
%% LETRAS MATHBB
%%
\newcommand{\R}{{\mathbb R}}
\newcommand{\Z}{{\mathbb Z}}
\newcommand{\C}{{\mathbb C}}
\newcommand{\Q}{{\mathbb Q}}
\newcommand{\N}{{\mathbb N}}
\newcommand{\E}{{\mathbb E}}
\newcommand{\V}{{\mathbb V}}
\renewcommand{\P}{{\mathbb P}}
\newcommand{\II}{{\mathbb I}}
\newcommand{\OO}{{\mathbb O}}
\newcommand{\Sb}{{\mathbb S}}
\newcommand{\Tb}{{\mathbb T}}

%%
%% LETRAS CALIFGRAFICAS
%%
\newcommand{\F}{{\mathcal F}}
\newcommand{\Hc}{{\mathcal H}}
\newcommand{\Tc}{{\mathcal T}}
\newcommand{\Ec}{{\mathcal E}}
\newcommand{\Vc}{{\mathcal V}}
\newcommand{\A}{{\mathcal A}}
\newcommand{\J}{{\mathcal J}}
\newcommand{\G}{{\mathcal G}}
\newcommand{\U}{{\mathcal U}}
\newcommand{\B}{{\mathcal B}}
\newcommand{\Sw}{{\mathcal S}}
\newcommand{\D}{{\mathcal D}}
\newcommand{\M}{{\mathcal M}}
\newcommand{\LL}{{\mathcal L}}

%%
%% LETRAS MATHBF
%%
\newcommand{\ab}{{\mathbf a}}
\newcommand{\bb}{{\mathbf b}}
\newcommand{\cc}{{\mathbf c}}
\newcommand{\ee}{{\mathbf e}}
\newcommand{\ff}{{\mathbf f}}
\newcommand{\gb}{{\mathbf g}}
\newcommand{\hh}{{\mathbf h}}
\newcommand{\n}{{\mathbf n}}
\newcommand{\mm}{{\mathbf m}}
\newcommand{\pp}{{\mathbf p}}
\newcommand{\qq}{{\mathbf q}}
\newcommand{\sbf}{{\mathbf s}}
\newcommand{\tb}{{\mathbf t}}
\newcommand{\uu}{{\mathbf u}}
\newcommand{\vv}{{\mathbf v}}
\newcommand{\ww}{{\mathbf w}}
\newcommand{\xx}{{\mathbf x}}
\newcommand{\yy}{{\mathbf y}}
\newcommand{\zz}{{\mathbf z}}
\newcommand{\I}{{\mathbf 1}}
\newcommand{\Pb}{{\mathbf P}}
\newcommand{\XX}{{\mathbf X}}
\newcommand{\YY}{{\mathbf Y}}

%%
%% LETRAS MATHFRAK
%%

\newcommand{\Gk}{{\mathfrak G}}


%%
%% TEOREMAS
%%
\newtheorem{teo}{Teorema}[section]
\newtheorem{lema}[teo]{Lema}
\newtheorem{prop}[teo]{Proposición}
\newtheorem{cor}[teo]{Corolario}

%%
%% REMARK
%%
\theoremstyle{remark}
\newtheorem{obs}[teo]{Observación}

%%
%% DEFINICION
%%
\theoremstyle{definition}
\newtheorem{defi}[teo]{Definición}
\newtheorem{ejem}[teo]{Ejemplo}
\newtheorem{ejer}[teo]{Ejercicio}
\newtheorem{nota}[teo]{Notación}
\newtheorem{ejercicio}{Ejercicio}[chapter]


%%
%% NUMERACION
%%
\numberwithin{subsection}{section}
\numberwithin{equation}{section}

%%
%% TITULO
%%


\vfuzz7pt \hfuzz7pt
\externaldocument[N-]{../notas/Notas-TO}
\newcommand{\encabezado}[1]{%
\noindent \includegraphics[scale=2.1]{logo-dm-wide.png} \hskip 9.5cm
\includegraphics[scale=0.08]{logo-exactas-uba.jpg}
\smallskip
\usefont{T1}{phv}{m}{n}
\begin{center}{\Large\sc Teoría de Transporte Óptimo}\end{center}
\bigskip
\noindent{\sc #1}\\[-0.4cm]
\hrule
\normalfont}
\newcommand{\cuadernolink}[2]{\noindent\textbf{Cuaderno:} \emph{#1} (\href{https://colab.research.google.com/github/jfbonder/transporte-optimo/blob/main/cuadernos/#2}{abrir en Colab}).\par\medskip}

\begin{document}
\renewcommand{\thechapter}{6}
\encabezado{Práctica 6: Potenciales de Kantorovich}
\bigskip
\pagestyle{fancy}
\fancyhead[L]{{\footnotesize{\usefont{T1}{phv}{m}{n}{\sc Teoría de Transporte Óptimo}}}}
\fancyhead[R]{\footnotesize \usefont{T1}{phv}{m}{n}{\sc Práctica 6}}
\thispagestyle{plain}
\noindent{\footnotesize Los ejercicios son los de la sección de ejercicios del Capítulo~6 de las notas del curso, con la misma numeración; las referencias a teoremas, proposiciones y secciones remiten a las notas (\url{https://github.com/jfbonder/transporte-optimo}).}
\medskip

\cuadernolink{Dualidad}{02-dualidad.ipynb}


\begin{ejercicio}[Propiedades de la $c$-transformada]\label{ejer.ctransf.propiedades}
Sea $c\colon X\times Y\to\R$ acotada y sea $\chi\colon X\to\R$ acotada.
\begin{enumerate}[label=(\alph*)]
\item Probar que $\chi^{c\bar c}\ge\chi$ y que $\chi^{c\bar cc}=\chi^c$.
Deducir que la operación $\chi\mapsto\chi^{c\bar c}$ es idempotente.
\item Recordemos (Definición~\ref{N-def.cconcava}) que $\varphi\colon X\to\R$
es $c$-cóncava si $\varphi=\zeta^{\bar c}$ para alguna
$\zeta\colon Y\to\R\cup\{-\infty\}$ no idénticamente $-\infty$. Probar que $\varphi$ es $c$-cóncava si y solo si $\varphi=\varphi^{c\bar c}$.
\item Probar que $(\chi+s)^c=\chi^c-s$ para $s\in\R$, y que
$\chi\le\tilde\chi$ implica $\chi^c\ge\tilde\chi^c$.
\item Probar que el ínfimo puntual de una familia de funciones
$c$-cóncavas es $c$-cóncava (o idénticamente $-\infty$).
\end{enumerate}
\end{ejercicio}

\begin{ejercicio}[Equicontinuidad del ínfimo]\label{ejer.equicontinuidad.inf}
Sea $(X,d)$ un espacio métrico, $A$ un conjunto y $\mathcal F$ una familia de
funciones $f\colon X\times A\to\R$ tal que
$\{f(\cdot,\alpha):f\in\mathcal F,\ \alpha\in A\}$ es equicontinua (con un
módulo de continuidad $\omega$ común). Probar que la familia
$\{x\mapsto\inf_{\alpha\in A}f(x,\alpha):f\in\mathcal F\}$ es equicontinua
con el mismo módulo $\omega$, siempre que los ínfimos sean finitos. Explicar
en qué punto de la demostración del Teorema~\ref{N-teo.existencia.potenciales}
se usa este hecho (Proposición~\ref{N-prop.ctransf.modulo}).
\end{ejercicio}

\begin{ejercicio}[Costos que son distancias]\label{ejer.ctransf.distancia}
Sea $X=Y$ y $c=d$ una distancia en $X$ (acotada, para simplificar).
\begin{enumerate}[label=(\alph*)]
\item Probar que $\varphi\colon X\to\R$ es $c$-cóncava si y solo si es
$1$-Lipschitz respecto de $d$, y que en ese caso $\varphi^c=-\varphi$.
\item Deducir que, para este costo, el problema dual se escribe
\[
\sup\Bigl\{\int\varphi\,d(\mu-\nu):\ \varphi\ \text{$1$-Lipschitz}\Bigr\}.
\]
(Este es el contenido del Lema~\ref{N-lema.KR.ctransformada} y del
Teorema~\ref{N-teo.KR} para $d(x,y)=|x-y|$.)
\item Aplicarlo a $c(x,y)=|x-y|^\alpha$ en $\R^d$ con $0<\alpha<1$: probar
que es una distancia y que
\[
\min_{\Pi(\mu,\nu)}\int|x-y|^\alpha\,d\pi
=\sup\Bigl\{\int u\,d(\mu-\nu):\ [u]_{0,\alpha}\le1\Bigr\},
\qquad
[u]_{0,\alpha}=\sup_{x\ne y}\frac{|u(x)-u(y)|}{|x-y|^\alpha},
\]
para $\mu,\nu$ con soporte compacto.
\end{enumerate}
\end{ejercicio}

\begin{ejercicio}[$c$-transformadas para el costo cuadrático]\label{ejer.ctransf.cuadratico}
Sea $c(x,y)=|x-y|^2$ en $\R^d$.
\begin{enumerate}[label=(\alph*)]
\item Para $\chi(x)=\lambda|x|^2$, calcular $\chi^c$ y mostrar que es finita
si y solo si $\lambda<1$; en ese caso $\chi^c(y)=-\frac{\lambda}{1-\lambda}|y|^2$.
\item Para $\chi(x)=a\cdot x+b$, calcular $\chi^c$ y $\chi^{c\bar c}$.
\item Probar que $\varphi$ es $c$-cóncava si y solo si
$u(x)=\frac12\bigl(|x|^2-\varphi(x)\bigr)$ es convexa y semicontinua
inferiormente (la normalización de la Sección~\ref{N-sec.costo.cuadratico}), y
probar que $\varphi^c(y)=|y|^2-2\bar u(y)$, donde $\bar u$ es la
transformada de Legendre de $u$.
\end{enumerate}
\end{ejercicio}

\begin{ejercicio}[La mejora de potenciales se estaciona en el caso discreto]\label{ejer.mejora.discreta}
En el problema discreto (Capítulo~\ref{N-cap.dual.discreto}), la
Proposición~\ref{N-prop:mejora-potenciales} dice que $(\varphi,\psi)\mapsto
(\varphi^{c\bar c},\varphi^c)$ no empeora el funcional dual.
\begin{enumerate}[label=(\alph*)]
\item Probar que, partiendo de cualquier $\varphi\in\R^N$, la iteración
$\varphi\mapsto\varphi^{c\bar c}$ se detiene en un paso: $\varphi^{c\bar c}$
ya es $c$-cóncava (Ejercicio~\ref{ejer.ctransf.propiedades}).
\item Mostrar con el ejemplo del Ejercicio~\ref{N-ejer.PL.2x2} que un par
$(\varphi^{c\bar c},\varphi^c)$ no tiene por qué ser óptimo: la mejora de
potenciales produce pares admisibles «ajustados» pero no resuelve el
problema dual. (En el Capítulo~\ref{N-cap.Sinkhorn} se verá que la versión
\emph{suavizada} de esta iteración sí converge al óptimo del problema
regularizado.)
\end{enumerate}
\end{ejercicio}

\subsection*{Ejercicios computacionales}

Los ejercicios siguientes se resuelven con el cuaderno \emph{Dualidad}, que
acompaña los Capítulos~\ref{N-cap.PL} a~\ref{N-cap.dual} (véase el
Apéndice~\ref{N-ap.cuadernos}); los enunciados figuran también al final del
cuaderno. Las referencias a secciones son a las del cuaderno.

\begin{ejercicio}[El $2\times2$ y su degeneración]\label{ejc.dual.2x2}
En el ejemplo $2\times2$ de la Sección~8 del cuaderno, reemplazar
$c_{11}=1$ por un parámetro $s$ y calcular el plan óptimo para $s\in[0,6]$.
Hallar el valor de $s$ en el que cambia el vértice óptimo y verificar que allí
todo el segmento $\Pi(a,b)$ es óptimo y los potenciales pierden unicidad (más
allá de la constante aditiva).
\end{ejercicio}

\begin{ejercicio}[Potenciales únicos y grafo conexo]\label{ejc.dual.conexo}
En el ejemplo $3\times3$ del Ejercicio~\ref{N-ejer.discreto.3x3} (Sección~9
del cuaderno), modificar las masas $a$ manteniendo $b$ hasta que el grafo de
soporte $G_\pi$ del plan óptimo sea conexo, y verificar que entonces $\psi_3$
queda determinado ($t_{\min}=t_{\max}$). Formular la conjetura general y
compararla con el Ejercicio~\ref{N-ejer.discreto.unicidad}.
\end{ejercicio}

\begin{ejercicio}[Holgura complementaria y saturación espuria]\label{ejc.dual.holgura}
En el ejemplo aleatorio de la Sección~1 del cuaderno, contar los pares $(i,j)$
con $\pi_{ij}>0$ y los pares con holgura $c_{ij}-\varphi_i-\psi_j<10^{-9}$; el
segundo conjunto contiene al primero (Ejercicio~\ref{N-ejer.PL.holgura}).
Repetir para $50$ instancias aleatorias y estimar con qué frecuencia hay
saturación sin masa. ¿Qué ocurre si los puntos $x_i$, $y_j$ se toman en una
grilla regular en lugar de al azar?
\end{ejercicio}

\begin{ejercicio}[Mejora de potenciales]\label{ejc.dual.mejora}
Sea $\mathcal D^*=\max\mathcal D$ el valor óptimo del problema dual discreto
(que coincide con el costo óptimo). Para $100$ vectores $\varphi$ aleatorios,
calcular la fracción de la brecha
dual $\mathcal D^*-\mathcal D(\varphi,\varphi^c)$ que cierra un paso de mejora
$\varphi\mapsto\varphi^{c\bar c}$ (Proposición~\ref{N-prop:mejora-potenciales})
y graficar su histograma. ¿Hay algún $\varphi$ inicial para el cual un solo
paso alcance el óptimo? (\emph{Sugerencia:} un paso alcanza el óptimo
exactamente cuando $\varphi^c$ ya es un $\psi$ óptimo. Partir de un
$\varphi^*$ óptimo y \emph{bajar} su valor en un índice $i$ que no realice
ningún mínimo $\min_i(c_{ij}-\varphi^*_i)$. Comprobar, en cambio, que un
$\varphi^*$ perturbado al azar no alcanza el óptimo ni con perturbaciones de
tamaño $10^{-6}$: la mejora no es continua.) Comparar con el
Ejercicio~\ref{ejer.mejora.discreta}.
\end{ejercicio}

\begin{ejercicio}[Costo $0$--$1$]\label{ejc.dual.01}
Verificar la fórmula $\min_{\Pi(a,b)}\langle C,\pi\rangle=\frac12\|a-b\|_1$
del Ejercicio~\ref{N-ejer.discreto.01} para $200$ pares $(a,b)$ aleatorios
sobre los mismos puntos. Tomar después soportes \emph{distintos} para $\mu$ y
$\nu$ (ningún $x_i$ coincide con un $y_j$) y observar que el costo óptimo es
siempre $1$: explicar por qué y qué relación tiene con la variación total de
$\mu-\nu$.
\end{ejercicio}

\end{document}
